% ==========================================================================
%  RSiena Post-Estimation Framework: Overview
%
%  Describes the unified pipeline for predict, marginalEffects (AME),
%  fitted, residuals, relative importance, entropy, and
%  selection/influence tables.
% ==========================================================================

\documentclass[11pt,a4paper]{article}
\usepackage[utf8]{inputenc}
\usepackage{amsmath,amssymb}
\usepackage{booktabs}
\usepackage{hyperref}
\usepackage{enumitem}
\usepackage[margin=2.5cm]{geometry}

\title{RSiena Post-Estimation Framework\\{\large Architecture \& User Guide}}
\author{Daniel Gotthardt}
\date{\today}

\begin{document}
\maketitle

% ------------------------------------------------------------------
\section{Overview}
% ------------------------------------------------------------------

The RSiena post-estimation framework provides a unified pipeline for
computing derived quantities from estimated Stochastic Actor-oriented
Models (SAOMs).  All quantities that go through the shared pipeline
have three layers:

\begin{enumerate}[leftmargin=*]
  \item \textbf{Change contributions} (wide struct) --- the matrix of
    objective function change statistics, one row per ego--choice (static)
    or ministep--choice (dynamic).
  \item \textbf{Per-$\theta$ prediction functions} (\texttt{predictFun}) ---
    compute the quantity of interest for a given parameter vector.
  \item \textbf{Uncertainty orchestration} (\texttt{sienaPostestimate}) ---
    draws $\theta^{(s)} \sim \mathcal{N}(\hat\theta, \hat\Sigma_\theta)$
    and aggregates uncertainty.
\end{enumerate}

Table~\ref{tab:inventory} gives a complete inventory of post-estimation
functions, their S3 method status, and whether they use the shared
uncertainty pipeline.

\begin{table}[ht]
\centering
\caption{Post-estimation function inventory (current state).}
\label{tab:inventory}
\small
\begin{tabular}{llccccl}
\toprule
Function & Class & \texttt{print} & \texttt{summary} & \texttt{plot}
  & Pipeline & Notes \\
\midrule
\texttt{predict.sienaFit}
  & (data.frame) & -- & -- & -- & Yes & \\
\texttt{marginalEffects}
  & (data.frame) & -- & -- & -- & Yes & \\
\texttt{interpret\_size}
  & \texttt{sienaRI} & \checkmark & \checkmark & \checkmark & Yes$^*$
  & $^*$optional; delegates to \texttt{relativeImportance} \\
\texttt{entropy}
  & (data.frame) & -- & -- & -- & Yes & \\
\texttt{selectionTable}
  & \texttt{selectionTable} & \checkmark & -- & -- & No & \\
\texttt{influenceTable}
  & \texttt{influenceTable} & \checkmark & -- & -- & No & \\
\bottomrule
\end{tabular}
\end{table}

Two dispatch generics exist: \texttt{marginalEffects} (dispatches on
\texttt{sienaFit} with full uncertainty, or \texttt{sienaEffects} for
point-estimate only) and \texttt{entropy} (dispatches on \texttt{sienaFit}
or \texttt{default}).  \texttt{interpret\_size} dispatches on
\texttt{sienaFit} or \texttt{sienaEffects} and delegates RI computation to
\texttt{relativeImportance}.  The functions \texttt{interpret\_selection} and
\texttt{interpret\_influence} are thin wrappers around \texttt{selectionTable}
and \texttt{influenceTable}, respectively.


% ------------------------------------------------------------------
\section{Quantities}
% ------------------------------------------------------------------

\subsection{Predicted Probabilities (\texttt{predict.sienaFit})}
For each ego--choice pair, the change probability is
\[
  P_j = \frac{\exp(u_j)}{\sum_h \exp(u_h)},
  \qquad u_j = \sum_k \theta_k \, \delta_k(j),
\]
where $\delta_k(j)$ are the change contributions for effect~$k$ and choice~$j$.
The \texttt{tieProb} variant adjusts for density direction:
$r_{ij} = P_j$ if $x_{ij}=0$ (creation), $r_{ij} = 1 - P_j$ if $x_{ij}=1$
(dissolution).

\subsection{Average Marginal Effects (\texttt{marginalEffects})}

The first difference uses one of two counterfactual update formulas,
selected by perturbation mode:

\paragraph{Alter (one-alternative) mode.}
When the change statistic varies across alters
(e.g.\ density, reciprocity):
\[
  P_j^{\mathrm{cf}}
    = \frac{P_j \exp(\Delta u_j)}{1 - P_j + P_j \exp(\Delta u_j)}.
\]

\paragraph{Ego-wide mode.}
When the change statistic is constant across alters
(e.g.\ \texttt{egoX}):
\[
  P_j^{\mathrm{cf}}
    = \frac{P_j \exp(\Delta u_j)}{\sum_k P_k \exp(\Delta u_k)}.
\]

Here $\Delta u_j = \theta_e \delta_e$ for a unit change in effect~$e$.
The update formula is implemented in C++ (\texttt{mlogit\_update}) and
wrapped by the R helper \texttt{mlogit\_update\_r()}, which accepts
string perturbation types (\texttt{"alter"}/\texttt{"ego"}).

Auto-detection is based on the \texttt{interactionType} column of the
effects object (\texttt{"ego"} $\to$ ego-wide, otherwise alter).  Users
can override via \texttt{perturbType1}/\texttt{perturbType2}.

Second differences compose two such counterfactuals to measure
interaction effects.

\paragraph{Mass contrasts.}
In ego-wide mode, the output includes two extra columns:
\texttt{massCreation} ($\Delta P_i^+$) and \texttt{massDissolution}
($\Delta P_i^-$), summing dyad-level change-probability first differences
within the creation and dissolution risk sets per ego.

\paragraph{Risk ratios.}
\texttt{mainEffect = "riskRatio"} returns $P'/P$ instead of $P'-P$. Not fully 
tested yet for different perturbation types and second differences, but 
the same auto-detection should apply.

\subsection{Relative Importance (\texttt{interpret\_size})}
For each ego, the $L^1$ distance between the full-$\theta$ distribution
and each leave-one-out distribution (effect~$k$ set to zero) measures
the importance of effect~$k$.  Normalising gives relative importance:
\[
  RI_k(i) = \frac{\sum_j |P_j - P_j^{(-k)}|}
                  {\sum_{k'} \sum_j |P_j - P_j^{(-k')}|}.
\]
When \texttt{uncertainty = TRUE}, the calculation now runs through
\texttt{sienaPostestimate}; the default path calls the legacy
\texttt{sienaRI()} directly.

The return object has class \texttt{sienaRI} with \texttt{print},
\texttt{summary}, and \texttt{plot} methods (base~R graphics:
\texttt{barplot} for static, \texttt{matplot} for dynamic).

\subsection{Entropy (\texttt{entropy})}
The normalised entropy (degree of certainty) is
\[
  R_H = 1 - \frac{H}{\log J}, \qquad
  H = -\sum_j p_j \log p_j,
\]
as defined in Snijders (2004) and now has a separate dispatch method. 
The \texttt{sienaFit} method computes $R_H$ for each ego--choice pair, 
while the \texttt{default} method computes $R_H$ for a single probability
vector (e.g.\ the average distribution across egos).  The former can be used to 
identify egos with more or less deterministic choice patterns; the latter can be used to summarise overall model fit (lower entropy indicates better fit).

\subsection{Selection and Influence Tables}

\texttt{selectionTable()} and \texttt{influenceTable()} compute
interpretive summaries of network--behaviour co-evolution models.
Both return typed data frames (\texttt{selectionTable},
\texttt{influenceTable}) with \texttt{print} methods but no
\texttt{plot} or \texttt{summary}.  They do not use the shared
uncertainty pipeline.



% ------------------------------------------------------------------
\section{Pipeline Architecture}
% ------------------------------------------------------------------

\begin{verbatim}
  Callers (predict.sienaFit / marginalEffects / entropy.sienaFit)
    build: specs (via makeSpec) + contribFun (via makeContribFun/chainStore)
                                     |
                  sienaPostestimate(specs, contribFun, thetaHat, covTheta, ...)
                                     |
                     +---------------+----------------+
                     |                                |
               point estimate                   uncertainty
                     |                                |
          makeEstimatorFun(specs,         makeEstimatorFun(specs,
            contribFun, ...)               uncertContribFun, ...)
                     |                                |
            estimatorFun(thetaHat)         drawSimBatch(estimatorFun,
               → named list                 thetaHat, covTheta, ...)
               of data.frames              → named list of stacked
                                             sim data.frames
                     |                                |
                     +---------------+----------------+
                                     |
                     aggregatePostEstimation(expects, raw_sims_list, specs)
                                     |
                          Final result (named list of data.frames)
\end{verbatim}

\noindent\textbf{contribFun} is a closure \texttt{function(theta)} $\to$
wide contrib struct.  For static models it ignores \texttt{theta} and
returns the precomputed struct.  For dynamic models it runs
\texttt{getDynamicChangeContributions} in batches via the \texttt{chainStore}
interface, accumulating partial estimates with streaming sum/count accumulators
(\texttt{makeEstimatorFun} handles the n3 batching internally, transparent to
\texttt{drawSimBatch}).

\noindent\textbf{specs} is a named list of per-effect spec entries (built
by \texttt{makeSpec}).  Each entry carries \texttt{predictFun},
\texttt{outcomeName}, \texttt{level}, \texttt{condition}, and effect-specific
arguments.  The multi-effect case (N~effects in one \texttt{marginalEffects}
call) runs all N~predictions inside a single closure invocation, sharing the
contrib struct and baseline probabilities.

% ------------------------------------------------------------------
\section{Key Helper Functions}
% ------------------------------------------------------------------

\begin{description}[style=nextline]
  \item[\texttt{makeSpec(predictFun, outcomeName, ...)}]
    Construct and validate a spec entry.  Required fields: \texttt{predictFun},
    \texttt{outcomeName}.  Optional: \texttt{level}, \texttt{condition},
    \texttt{aggFun}, and caller-specific extras (e.g.\ \texttt{massContrasts}).
  \item[\texttt{makeContribFun(mode, store, effects, ...)}]
    Wraps a \texttt{chainStore} backend into a \texttt{function(theta)}
    $\to$ wide contrib struct closure.  Handles
    \texttt{"memory"}/\texttt{"disk"}/\texttt{"simulate"} modes and calls
    \texttt{flattenAndEnrichWide} internally.
  \item[\texttt{chainStore\_disk(chains, batchSize, dir)}]
    Flush raw chains to RDS files before forking workers; returns a
    \texttt{chainStore} with \texttt{getBatch}/\texttt{cleanup} interface.
    Zero copy-on-write overhead under \texttt{FORK} parallelism.
  \item[\texttt{chainStore\_simulate(dynArgs, batchSize, n3Total)}]
    On-the-fly backend: re-runs \texttt{getDynamicChangeContributions} for
    each batch rather than reading stored chains.  Used when
    \texttt{useChangeContributions = FALSE}.
  \item[\texttt{makeEstimatorFun(specs, contribFun, nBatches, type, ...)}]
    Core closure builder.  Returns \texttt{function(theta)} $\to$ named list
    of data.frames (one per spec).  Internally loops over chain batches,
    accumulates partial sum/count accumulators, and finalises the batch mean.
    n3~batching is fully encapsulated here.
  \item[\texttt{drawSimBatch(estimatorFun, thetaHat, covTheta, nsim, ...)}]
    Theta-draw loop: samples $\theta^{(s)} \sim \mathcal{N}(\hat\theta,
    \hat\Sigma)$, calls \texttt{estimatorFun} per draw.  Supports batch-to-disk
    persistence, resume, PSOCK/FORK parallelism via
    \texttt{openCluster}/\texttt{closeCluster}.
  \item[\texttt{aggregatePostEstimation(expects, raw\_sims\_list, specs, ...)}]
    Combines the point-estimate named list (\texttt{expects}) with the
    stacked sim draws (\texttt{raw\_sims\_list}) into final merged output.
    Calls \texttt{makeUncertaintySummarizer} and \texttt{mergeEstimates}.
  \item[\texttt{agg(outcomeName, data, level, condition, sum\_fun)}]
    Flexible group-by aggregation (data.table fast path + base~R fallback).
    \texttt{aggMulti} handles the N-effect case.
  \item[\texttt{makeUncertaintySummarizer(...)}]
    Builds summary function (mean, SD, CI, MCSE) from user-specified
    probability arguments.
  \item[\texttt{alignThetaNoRate(theta, effectNames, ans)}]
    Aligns $\theta$ to effect names, dropping rate effects.
  \item[\texttt{mergeEstimates(pointEst, uncertainty)}]
    Joins point estimate with uncertainty summary.
  \item[\texttt{openCluster / closeCluster}]
    PSOCK/FORK cluster lifecycle.  \texttt{openCluster} exports the
    \texttt{estimatorFun} closure (which captures all required state)
    to workers.
  \item[\texttt{planBatch(data, depvar, nsim, nbrNodes, ...)}]
    Compute \texttt{batchSize} from data dimensions, memory budget, and
    parallel topology.  Formula: peak $= W \cdot U_\text{call} +
    B_\text{batch} \cdot U_\text{agg}$; solves for the largest batch that
    fits within the budget.
\end{description}


% ------------------------------------------------------------------
\section{Limitation: Network DVs Only}
% ------------------------------------------------------------------

The shared pipeline currently supports only \textbf{network dependent
variables}.  Behaviour DVs (discrete and continuous) trigger a
\texttt{stop()} in \texttt{marginalEffects}.

Continuous behaviour DVs in RSiena use an SDE model (Ornstein--Uhlenbeck
with a Wiener process) rather than the multinomial-logit choice model.
Marginal effects for continuous behaviour would require a different
perturbation approach (shifting drift parameters rather than choice
utilities).


% ------------------------------------------------------------------
\section{Backward Compatibility}
% ------------------------------------------------------------------

The original \texttt{sienaRI}, \texttt{sienaRIDynamics},
\texttt{print.sienaRI}, \texttt{summary.sienaRI}, \texttt{plot.sienaRI}
remain fully functional.  The default path in \texttt{interpret\_size.sienaFit}
still calls \texttt{sienaRI()} directly; the new path is activated by
\texttt{uncertainty = TRUE}.


% ------------------------------------------------------------------
\section{Outlook}
% ------------------------------------------------------------------

The near-term priority is adding an S3 class (\texttt{"sienaMarginal"})
to the \texttt{marginalEffects} return value, together with
\texttt{print}, \texttt{summary}, and \texttt{plot} methods (base~R
only, following \texttt{plot.sienaRI}).  This would bring
\texttt{marginalEffects} in line with the other postestimation tools
that already carry typed classes.

The \texttt{drawSimBatch} function is the current draw loop; a planned
Step~7 (see \texttt{docs/design/postestimation/sienaPostEstimationNewPlan.md})
will merge \texttt{drawSimBatch} into a unified \texttt{drawSim} that also
supports list-valued estimators.

Longer-term extensions include support for behaviour dependent variables
(discrete and continuous) and period-level expected change counts that
integrate the rate and choice components.


% ------------------------------------------------------------------
\section{References}
% ------------------------------------------------------------------

\begin{itemize}
  \item Gotthardt, D.\ and Steglich, C.E.G.\ (2026).
    Marginal effects for the stochastic actor-oriented model.
  \item Indlekofer, N.\ and Brandes, U.\ (2013).
    Relative Importance of Effects in Stochastic Actor-oriented Models.
    \emph{Network Science}, 1, 278--304.
  \item Snijders, T.A.B.\ (2004).
    Explained Variation in Dynamic Network Models.
    \emph{Math\'{e}matiques et Sciences Humaines}, 168, 31--41.
\end{itemize}

\end{document}
